\documentclass[10pt,a4paper]{amsart}
\usepackage[margin=19mm]{geometry}
\usepackage{amsmath,amssymb,mathtools}
\usepackage[scr=rsfs,cal=euler]{mathalfa}
\usepackage{xfrac}
\usepackage[unicode,hidelinks,bookmarks=false]{hyperref}
\newcommand{\ie}{i.e.\ }
\newcommand{\cf}{cf.\ }
\newcommand{\resp}{respectively\ }
\newcommand{\Subsection}{\subsection}
\providecommand{\nobreakdash}{\nobreak\hbox{-}}
\setlength{\emergencystretch}{2em}
\multlinegap0pt
\usepackage{amssymb}
\newtheorem{example}{Example}
\newtheorem{theorem}{Theorem}
\title{Quasi sdf-absorbing ideals in commutative rings}
\author{Violeta Leoreanu-Fotea, Ece Yetkin Celikel, Tarik Arabaci, Unsal Tekir}
\date{Source: arXiv:2605.05279v1 (CC BY 4.0)}
\begin{document}
\maketitle

\noindent\emph{Source and licence.} Quasi sdf-absorbing ideals in commutative rings. Version arXiv:2605.05279v1 (CC BY 4.0), distributed by arXiv under the Creative Commons Attribution 4.0 International licence, http://creativecommons.org/licenses/by/4.0/, as recorded in the arXiv licence metadata for this exact version and shown on the abstract page https://arxiv.org/abs/2605.05279v1. Full licence notice: \texttt{licenses/arxiv-2605.05279-CC-BY-4.0.txt}.\par\smallskip

\noindent\textit{Editorial scope.} Counted unit: the article's theorem theorem (source0.tex lines 755--758). The article's complete original proof of it is delivered with it (source0.tex lines 760--796, one \texttt{proof} environment of 37 lines). No mathematics is written, repaired or re-derived: the statement and the proof are the article's own lines, in the article's own notation, and no retained line is substituted or re-flowed.  Retained uncounted as the source-local context the proof needs: lines 584--586.  Nothing was omitted from any retained span, so the package carries no omission record.  No retained line contains a citation, so the wrapper carries no bibliography and no bibliography entry.  No retained line contains a figure environment, an image-inclusion command or drawing code, so no image is delivered and the package declares no asset dependency.  Editorial formatting only, in addition: an AMS \texttt{amsart} preamble that restates the article's own declarations and macros used by the retained lines, namely the environments \texttt{example}, \texttt{theorem} on separate counters, together with the standard packages those lines need.  The retained environments print in this document as Example 1, Theorem 1 in order of appearance, whereas the article prints the same items with the numbering of its own full document.  The complete native source of the article as distributed in the arXiv e-print for this exact version is bundled unchanged as \texttt{sources/199940/source0.tex}.
\begin{example}
\label{estar} \ 
\end{example}

\begin{theorem}
Let $K$ be a field. The polynomial ring $K[x]$ always satisfies Condition
$(\ast)$.
\end{theorem}

\begin{proof}
If $char(K)=2,$ then the claim is clear by Example \ref{estar}(2). So, suppose
that $char(K)\neq2$. Let $I\subseteq K[x]$ be a quasi sdf-absorbing ideal. By
definition, its radical $L=\sqrt{I}$ is an sdf-absorbing ideal, meaning that
for any $a,b\in K[x]$, if $a^{2}-b^{2}\in L$, then $(a-b)\in L$ or $(a+b)\in
L$. We aim to show that $I$ is a primary ideal. \newline Since $K[x]$ is a
Principal Ideal Domain (PID), every ideal is generated by a single polynomial.
Let $I=(f)$ and $L=\sqrt{I}=(p)$. In a PID, the radical of an ideal generated
by $f$ is generated by the product of the distinct irreducible factors of $f$.
Thus, we can write $p=p_{1}p_{2}\dots p_{n},$ where $p_{i}$ are distinct
irreducible polynomials in $K[x]$. We aim to show $n=1$. Suppose $n>1$. We use
the Chinese Remainder Theorem to construct a polynomial $g$ such that:%
\[%
\begin{cases}
g\equiv1\pmod{p_1}\\
g\equiv-1\pmod{p_2 \dots p_n}
\end{cases}
\]
This is possible because $p_{1}$ and the product $p_{2}\dots p_{n}$ are
relatively prime. Let $a=g$ and $b=1$. Then $a^{2}-b^{2}=(g-1)(g+1)$. Since
$g\equiv1\pmod{p_1}$, we have $p_{1}\mid(g-1)$. Also, since $g\equiv
-1\pmod{p_2 \dots p_n}$, we have $(p_{2}\dots p_{n})\mid(g+1)$. Now, since
$p_{1}$ and $p_{2}\dots p_{n}$ are coprime, then $p\mid(g-1)(g+1)$ which means
$g^{2}-1\in(p)=L.$ Since $L$ is sdf-absorbing, either $(g-1)\in L$ or
$(g+1)\in L$. Hence, we have the following cases:

Case I: If $(g-1)\in L=(p_{1}\dots p_{n})$, then $p_{2}\mid(g-1)$. But our
construction says $g\equiv-1\pmod{p_2}$, so $g-1\equiv-2\pmod{p_2}$. As
$\text{char}(K)\neq2$, this implies $p_{2}$ is a unit, a contradiction.

Case II. If $(g+1)\in L$, then $p_{1}\mid(g+1)$. But $g\equiv1\pmod{p_1}$, so
$g+1\equiv2\pmod{p_1}$, again implying $p_{1}$ is a unit, a contradiction.

The assumption $n>1$ leads to a contradiction. Thus $n=1$, meaning $L=(p_{1})$
is prime. This ensures $I$ is a primary ideal, and specifically, an
sdf-absorbing primary ideal.
\end{proof}

\end{document}
